% BHU PYQ -- Manually transcribed & error-corrected
\documentclass[11pt,a4paper]{article}

\usepackage[a4paper,top=2.5cm,bottom=2.5cm,left=2.8cm,right=2.8cm,headheight=14pt]{geometry}
\usepackage{amsmath,amssymb}
\usepackage[T1]{fontenc}
\usepackage[utf8]{inputenc}
\usepackage{lmodern}
\usepackage{microtype}
\usepackage{fancyhdr}
\usepackage{enumitem}
\usepackage{hyperref}
\usepackage{lastpage}
\usepackage{parskip}

\hypersetup{colorlinks=true,urlcolor=black,linkcolor=black,
  pdftitle={BPT-603: Elements of Nuclear Physics -- BHU PYQ},pdfauthor={Banaras Hindu University}}

\pagestyle{fancy}\fancyhf{}
\renewcommand{\headrulewidth}{0.4pt}\renewcommand{\footrulewidth}{0.4pt}
\fancyhead[L]{\small   \textit{Banaras Hindu University}}
\fancyhead[R]{\small   \textit{BPT-603: Elements of Nuclear Physics}}
\fancyfoot[C]{\small Page  \thepage\ of \pageref{LastPage}}


\newlist{parts}{enumerate}{2}
\setlist[parts,1]{label=(\alph*),leftmargin=*,itemsep=5pt,topsep=3pt}
\setlist[parts,2]{label=(\roman*),leftmargin=*,itemsep=3pt,topsep=2pt}

\newcommand{\pts}[1]{\hfill   \textit{[#1]}}


\begin{document}


\begin{center}
  {\large\bfseries BANARAS HINDU UNIVERSITY}\\[2pt]
  {
\normalsize B.Sc. (Hons.) Semester VI Examination, 2022-23}\\[6pt]
  \rule{0.8\linewidth}{0.5pt}\\[6pt]
  {\large\bfseries Paper No. BPT-603: Elements of Nuclear Physics}\\[4pt]
  {
\normalsize Time: 3 Hours\hfill Maximum Marks: 70}
\end{center}
\rule{\linewidth}{0.4pt}
\smallskip



\noindent   \textit{Instructions: Answer all questions. Symbols have their usual meanings.}
\bigskip




\noindent   \textbf{Question 1.} Answer any   \textbf{five} questions as given below: \pts{5 $ \times$ 4 = 20}

\begin{parts}
  \item All nuclides are made up of a neutron-proton mixture that we can call nuclear matter. What is its density?
  \item Which of the following reactions are possible and why?
    
\begin{parts}
      \item $n + n  o \Lambda^0 + K^+$
      \item $p + p  o K^+ + \Sigma^+$
    \end{parts}
  \item How much energy is necessary to split up an alpha particle into its constituent nucleons? \\
    (Given: $M_n = 1.008665\, \text{u}$, $M_p = 1.007825\, \text{u}$, $M(^{4} \text{He}) = 4.00260\, \text{u}$).
  \item Discuss the $Q$-value formula for all three kinds of Beta decay.
  \item Carbon-11 transforms by $\beta^+$ decay. Determine the products of the reaction and the energy released. Ignore the mass of the emitted neutrino. (Given: mass of $^{11} \text{C} = 11.011434\, \text{u}$, mass of $^{11} \text{B} = 11.009305\, \text{u}$).
  \item A relic is found to give an activity count of $12\, \text{c.p.m.}$ (counts per minute) for each gram of carbon. If living trees give a count of $16\, \text{c.p.m.}$, find the approximate age of the relic. (Take the half-life of Carbon-14 as $5730$ years).
\end{parts}

\medskip\rule{\linewidth}{0.2pt}




\noindent   \textbf{Question 2.}

\begin{parts}
  \item Discuss experimental methods for the charge radius measurement of nuclei. Explain with a proper diagram. Also estimate the minimum energy of the incident ion for determining the radius of the target nucleus. \pts{7.5}
  \item Explain the need for introducing the colour degree of freedom for the quarks. \pts{5}
\end{parts}

\medskip\rule{\linewidth}{0.2pt}




\noindent   \textbf{Question 3.}

\begin{parts}
  \item Discuss the different properties of the nuclear force in detail along with supporting experimental evidence. \pts{7}
  \item The decay of $^{253} \text{Es}$ ($I = 7/2$, $\pi = +$) leads to a sequence of negative-parity states in $^{249} \text{Bk}$ with $I = 3/2, 5/2, 7/2, 9/2, 11/2, 13/2$. For each state, find the permitted values of the orbital angular momentum $l$ of the emitted alpha particle. \pts{5.5}
\end{parts}
\hfill   \textbf{OR}

\begin{parts}
  \item Explain the working of the Van de Graaff (VdG) accelerator using a suitable diagram. \pts{7.5}
  \item Describe the expression for the Coulomb potential barrier in Alpha decay. Calculate the Coulomb potential barrier experienced by an Alpha particle in the Alpha decay of $^{226} \text{Ra}$ ($Z = 88$). \pts{5}
\end{parts}

\medskip\rule{\linewidth}{0.2pt}




\noindent   \textbf{Question 4.}

\begin{parts}
  \item An even-$Z$, even-$N$ nucleus has the following sequence of levels above its $0^+$ ground state: $2^+$ ($89\, \text{keV}$), $4^+$ ($288\, \text{keV}$), $6^+$ ($585\, \text{keV}$), $0^+$ ($1050\, \text{keV}$), $2^+$ ($1129\, \text{keV}$). Draw an energy level diagram and show all probable $\gamma$ transitions and their dominant multipole assignments. \pts{6.5}
  \item For the following $\gamma$-transitions, give all permitted multipoles and indicate which multipole might be the most intense in the emitted radiation:
    
\begin{parts}
      \item $9/2^+  o 7/2^+$
      \item $1/2^-  o 7/2^-$
      \item $1^-  o 2^+$
      \item $4^+  o 2^+$
      \item $11/2^-  o 3/2^+$
      \item $3^+  o 3^+$
    \end{parts} \pts{6}
\end{parts}
\hfill   \textbf{OR}

\begin{parts}
  \item A beam of $8\, \text{MeV}$ $\alpha$-particles is incident upon a thin gold screen ($Z = 79$). Calculate the impact parameter for scattering an incident $\alpha$-particle through an angle of $110^\circ$. \pts{4.5}
  \item Give the quark constituents of the members of the meson $0^-$ octet, baryon $1/2^+$ octet, and baryon $3/2^+$ decuplet families. \pts{8}
\end{parts}

\medskip\rule{\linewidth}{0.2pt}




\noindent   \textbf{Question 5.}

\begin{parts}
  \item Discuss the photomultiplier tube (PMT) using a suitable diagram. Also explain gain and quantum efficiency. \pts{4.5}
  \item In the semi-empirical mass formula, what experimental evidence does the pairing term $\delta$ account for? For a given value of the mass number $A$, nuclear masses can be expressed as a quadratic function in $Z$. Using the minimum isobar formula, determine whether $^{142} \text{Xe}$ ($Z = 54$) is $\beta^+$ or $\beta^-$ unstable. \pts{8}
\end{parts}

\vfill
\rule{\linewidth}{0.4pt}

\begin{center}\small
  \href{https://bhu-exam.vercel.app/}{Click here to visit our website}\\[4pt]
  \href{https://me-aryan.vercel.app/}{Click here to visit our contributor}\\[4pt]
  \href{https://quashan-me.vercel.app/}{Click here to visit our contributor}
\end{center}

\end{document}
